% Options for packages loaded elsewhere
% Options for packages loaded elsewhere
\PassOptionsToPackage{unicode}{hyperref}
\PassOptionsToPackage{hyphens}{url}
\PassOptionsToPackage{dvipsnames,svgnames,x11names}{xcolor}
%
\documentclass[
  letterpaper,
  DIV=11,
  numbers=noendperiod]{scrartcl}
\usepackage{xcolor}
\usepackage[margin=1in]{geometry}
\usepackage{amsmath,amssymb}
\setcounter{secnumdepth}{-\maxdimen} % remove section numbering
\usepackage{iftex}
\ifPDFTeX
  \usepackage[T1]{fontenc}
  \usepackage[utf8]{inputenc}
  \usepackage{textcomp} % provide euro and other symbols
\else % if luatex or xetex
  \usepackage{unicode-math} % this also loads fontspec
  \defaultfontfeatures{Scale=MatchLowercase}
  \defaultfontfeatures[\rmfamily]{Ligatures=TeX,Scale=1}
\fi
\usepackage{lmodern}
\ifPDFTeX\else
  % xetex/luatex font selection
\fi
% Use upquote if available, for straight quotes in verbatim environments
\IfFileExists{upquote.sty}{\usepackage{upquote}}{}
\IfFileExists{microtype.sty}{% use microtype if available
  \usepackage[]{microtype}
  \UseMicrotypeSet[protrusion]{basicmath} % disable protrusion for tt fonts
}{}
\makeatletter
\@ifundefined{KOMAClassName}{% if non-KOMA class
  \IfFileExists{parskip.sty}{%
    \usepackage{parskip}
  }{% else
    \setlength{\parindent}{0pt}
    \setlength{\parskip}{6pt plus 2pt minus 1pt}}
}{% if KOMA class
  \KOMAoptions{parskip=half}}
\makeatother
% Make \paragraph and \subparagraph free-standing
\makeatletter
\ifx\paragraph\undefined\else
  \let\oldparagraph\paragraph
  \renewcommand{\paragraph}{
    \@ifstar
      \xxxParagraphStar
      \xxxParagraphNoStar
  }
  \newcommand{\xxxParagraphStar}[1]{\oldparagraph*{#1}\mbox{}}
  \newcommand{\xxxParagraphNoStar}[1]{\oldparagraph{#1}\mbox{}}
\fi
\ifx\subparagraph\undefined\else
  \let\oldsubparagraph\subparagraph
  \renewcommand{\subparagraph}{
    \@ifstar
      \xxxSubParagraphStar
      \xxxSubParagraphNoStar
  }
  \newcommand{\xxxSubParagraphStar}[1]{\oldsubparagraph*{#1}\mbox{}}
  \newcommand{\xxxSubParagraphNoStar}[1]{\oldsubparagraph{#1}\mbox{}}
\fi
\makeatother


\usepackage{longtable,booktabs,array}
\usepackage{calc} % for calculating minipage widths
% Correct order of tables after \paragraph or \subparagraph
\usepackage{etoolbox}
\makeatletter
\patchcmd\longtable{\par}{\if@noskipsec\mbox{}\fi\par}{}{}
\makeatother
% Allow footnotes in longtable head/foot
\IfFileExists{footnotehyper.sty}{\usepackage{footnotehyper}}{\usepackage{footnote}}
\makesavenoteenv{longtable}
\usepackage{graphicx}
\makeatletter
\newsavebox\pandoc@box
\newcommand*\pandocbounded[1]{% scales image to fit in text height/width
  \sbox\pandoc@box{#1}%
  \Gscale@div\@tempa{\textheight}{\dimexpr\ht\pandoc@box+\dp\pandoc@box\relax}%
  \Gscale@div\@tempb{\linewidth}{\wd\pandoc@box}%
  \ifdim\@tempb\p@<\@tempa\p@\let\@tempa\@tempb\fi% select the smaller of both
  \ifdim\@tempa\p@<\p@\scalebox{\@tempa}{\usebox\pandoc@box}%
  \else\usebox{\pandoc@box}%
  \fi%
}
% Set default figure placement to htbp
\def\fps@figure{htbp}
\makeatother





\setlength{\emergencystretch}{3em} % prevent overfull lines

\providecommand{\tightlist}{%
  \setlength{\itemsep}{0pt}\setlength{\parskip}{0pt}}



 


\KOMAoption{captions}{tableheading}
\makeatletter
\@ifpackageloaded{tcolorbox}{}{\usepackage[skins,breakable]{tcolorbox}}
\@ifpackageloaded{fontawesome5}{}{\usepackage{fontawesome5}}
\definecolor{quarto-callout-color}{HTML}{909090}
\definecolor{quarto-callout-note-color}{HTML}{0758E5}
\definecolor{quarto-callout-important-color}{HTML}{CC1914}
\definecolor{quarto-callout-warning-color}{HTML}{EB9113}
\definecolor{quarto-callout-tip-color}{HTML}{00A047}
\definecolor{quarto-callout-caution-color}{HTML}{FC5300}
\definecolor{quarto-callout-color-frame}{HTML}{acacac}
\definecolor{quarto-callout-note-color-frame}{HTML}{4582ec}
\definecolor{quarto-callout-important-color-frame}{HTML}{d9534f}
\definecolor{quarto-callout-warning-color-frame}{HTML}{f0ad4e}
\definecolor{quarto-callout-tip-color-frame}{HTML}{02b875}
\definecolor{quarto-callout-caution-color-frame}{HTML}{fd7e14}
\makeatother
\makeatletter
\@ifpackageloaded{caption}{}{\usepackage{caption}}
\AtBeginDocument{%
\ifdefined\contentsname
  \renewcommand*\contentsname{Table of contents}
\else
  \newcommand\contentsname{Table of contents}
\fi
\ifdefined\listfigurename
  \renewcommand*\listfigurename{List of Figures}
\else
  \newcommand\listfigurename{List of Figures}
\fi
\ifdefined\listtablename
  \renewcommand*\listtablename{List of Tables}
\else
  \newcommand\listtablename{List of Tables}
\fi
\ifdefined\figurename
  \renewcommand*\figurename{Figure}
\else
  \newcommand\figurename{Figure}
\fi
\ifdefined\tablename
  \renewcommand*\tablename{Table}
\else
  \newcommand\tablename{Table}
\fi
}
\@ifpackageloaded{float}{}{\usepackage{float}}
\floatstyle{ruled}
\@ifundefined{c@chapter}{\newfloat{codelisting}{h}{lop}}{\newfloat{codelisting}{h}{lop}[chapter]}
\floatname{codelisting}{Listing}
\newcommand*\listoflistings{\listof{codelisting}{List of Listings}}
\makeatother
\makeatletter
\makeatother
\makeatletter
\@ifpackageloaded{caption}{}{\usepackage{caption}}
\@ifpackageloaded{subcaption}{}{\usepackage{subcaption}}
\makeatother
\usepackage{bookmark}
\IfFileExists{xurl.sty}{\usepackage{xurl}}{} % add URL line breaks if available
\urlstyle{same}
\hypersetup{
  pdftitle={Unit 2: Ploidy --- How Many Copies?},
  colorlinks=true,
  linkcolor={blue},
  filecolor={Maroon},
  citecolor={Blue},
  urlcolor={Blue},
  pdfcreator={LaTeX via pandoc}}


\title{Unit 2: Ploidy --- How Many Copies?}
\author{}
\date{}
\begin{document}
\maketitle

\renewcommand*\contentsname{Table of contents}
{
\hypersetup{linkcolor=}
\setcounter{tocdepth}{3}
\tableofcontents
}

\subsection{What you're learning here}\label{what-youre-learning-here}

By the end, you should be able to:

\begin{itemize}
\tightlist
\item
  Define \textbf{ploidy} using our course definition
\item
  Tell the difference between a \textbf{diploid/haploid cell} and
  \textbf{diploid/haploid number}
\item
  Identify \textbf{homologous chromosomes} vs \textbf{sister chromatids}
\item
  Determine ploidy level (haploid, diploid, triploid, tetraploid, etc.)
  from a diagram
\end{itemize}

\begin{center}\rule{0.5\linewidth}{0.5pt}\end{center}

\subsection{The core definition (use this all
semester)}\label{the-core-definition-use-this-all-semester}

\begin{tcolorbox}[enhanced jigsaw, titlerule=0mm, coltitle=black, bottomrule=.15mm, colframe=quarto-callout-important-color-frame, title=\textcolor{quarto-callout-important-color}{\faExclamation}\hspace{0.5em}{Ploidy (course definition)}, arc=.35mm, opacityback=0, opacitybacktitle=0.6, colback=white, breakable, toprule=.15mm, rightrule=.15mm, bottomtitle=1mm, left=2mm, toptitle=1mm, colbacktitle=quarto-callout-important-color!10!white, leftrule=.75mm]

\textbf{Ploidy = the number of chromosomes in a complete set of
homologs.}

\end{tcolorbox}

A \emph{complete set} means: ``one of each chromosome type.''

\includegraphics[width=1\linewidth,height=\textheight,keepaspectratio]{../images/ploidy/ploidy_definition_set.png}
\emph{Figure 1. One complete set of homologs contains one of each
chromosome type (e.g., chromosome 1, chromosome 2, etc.).}

\begin{center}\rule{0.5\linewidth}{0.5pt}\end{center}

\subsection{Step 1: ``Diploid cell'' vs ``diploid number'' (they are not
the
same)}\label{step-1-diploid-cell-vs-diploid-number-they-are-not-the-same}

\textbf{Diploid cell} = has \textbf{two} versions of each chromosome
(two homologs). :contentReference{oaicite:2}\\
\textbf{Diploid number} = the \textbf{total} number of chromosomes in
that diploid cell.

\begin{tcolorbox}[enhanced jigsaw, titlerule=0mm, coltitle=black, bottomrule=.15mm, colframe=quarto-callout-note-color-frame, title=\textcolor{quarto-callout-note-color}{\faInfo}\hspace{0.5em}{Two symbols you must keep separate}, arc=.35mm, opacityback=0, opacitybacktitle=0.6, colback=white, breakable, toprule=.15mm, rightrule=.15mm, bottomtitle=1mm, left=2mm, toptitle=1mm, colbacktitle=quarto-callout-note-color!10!white, leftrule=.75mm]

\begin{itemize}
\tightlist
\item
  \textbf{2n} describes a \emph{diploid cell} (two homologs per
  chromosome type)
\item
  \textbf{2n} also helps you calculate the \emph{diploid number} once
  you know \textbf{n}
\end{itemize}

\end{tcolorbox}

\includegraphics[width=0.9\linewidth,height=\textheight,keepaspectratio]{../images/ploidy/diploid_cell_cartoon.png}
\emph{Figure 2. A diploid cell: two homologs for each chromosome type.}

\begin{center}\rule{0.5\linewidth}{0.5pt}\end{center}

\subsection{\texorpdfstring{Step 2: What does \textbf{n}
mean?}{Step 2: What does n mean?}}\label{step-2-what-does-n-mean}

\textbf{n} = the number of \textbf{distinct chromosome types} (the size
of one complete set). :contentReference{oaicite:3}

\includegraphics[width=0.95\linewidth,height=\textheight,keepaspectratio]{../images/ploidy/n_equals_unique_types.png}
\emph{Figure 3. n is the number of unique chromosome types in one set.}

\subsubsection{Example (small on
purpose)}\label{example-small-on-purpose}

If a pretend organism has two chromosome types (chromosome 1 and 2),
then: - \textbf{n = 2} - diploid cell is \textbf{2n} - diploid number is
\textbf{2 × 2 = 4}

\includegraphics[width=0.9\linewidth,height=\textheight,keepaspectratio]{../images/ploidy/2n_equals_4_example.png}
\emph{Figure 4. If n=2, then 2n=4 total chromosomes in a diploid cell.}

\begin{center}\rule{0.5\linewidth}{0.5pt}\end{center}

\subsection{Haploid: one complete set}\label{haploid-one-complete-set}

A \textbf{haploid cell} has \textbf{one} version of each chromosome type
(one complete set). :contentReference{oaicite:4}

\includegraphics[width=0.75\linewidth,height=\textheight,keepaspectratio]{../images/ploidy/haploid_cell_cartoon.png}
\emph{Figure 5. A haploid cell contains one complete set.}

\begin{tcolorbox}[enhanced jigsaw, titlerule=0mm, coltitle=black, bottomrule=.15mm, colframe=quarto-callout-important-color-frame, title=\textcolor{quarto-callout-important-color}{\faExclamation}\hspace{0.5em}{The cleanest shortcut}, arc=.35mm, opacityback=0, opacitybacktitle=0.6, colback=white, breakable, toprule=.15mm, rightrule=.15mm, bottomtitle=1mm, left=2mm, toptitle=1mm, colbacktitle=quarto-callout-important-color!10!white, leftrule=.75mm]

The \textbf{haploid number} is just \textbf{n}.

\end{tcolorbox}

\begin{center}\rule{0.5\linewidth}{0.5pt}\end{center}

\subsection{Homologous chromosomes: same genes, possibly different
alleles}\label{homologous-chromosomes-same-genes-possibly-different-alleles}

In a diploid cell, the two versions of the same chromosome are
\textbf{homologous chromosomes}. :contentReference{oaicite:5}

\textbf{What to notice} - same genes in the same order - alleles can
differ

\includegraphics[width=1\linewidth,height=\textheight,keepaspectratio]{../images/ploidy/homologs_same_genes.png}
\emph{Figure 6. Homologous chromosomes carry the same genes in the same
locations (loci).}

\begin{center}\rule{0.5\linewidth}{0.5pt}\end{center}

\subsection{Why chromosomes sometimes look like ``X''
shapes}\label{why-chromosomes-sometimes-look-like-x-shapes}

Before mitosis, the cell replicates DNA in \textbf{S-phase}.
:contentReference{oaicite:6}

After replication: - each chromosome has \textbf{two identical copies
stuck together} - those identical copies are \textbf{sister chromatids}
- they are joined at the \textbf{centromere}
:contentReference{oaicite:7}

\includegraphics[width=1\linewidth,height=\textheight,keepaspectratio]{../images/ploidy/s_phase_replication.png}
\emph{Figure 7. DNA replication in S-phase produces sister chromatids.}

\includegraphics[width=1\linewidth,height=\textheight,keepaspectratio]{../images/ploidy/sister_chromatids_vs_homologs.png}
\emph{Figure 8. Homologous chromosomes vs sister chromatids (do not mix
these up).}

\begin{tcolorbox}[enhanced jigsaw, titlerule=0mm, coltitle=black, bottomrule=.15mm, colframe=quarto-callout-note-color-frame, title=\textcolor{quarto-callout-note-color}{\faInfo}\hspace{0.5em}{Quick check}, arc=.35mm, opacityback=0, opacitybacktitle=0.6, colback=white, breakable, toprule=.15mm, rightrule=.15mm, bottomtitle=1mm, left=2mm, toptitle=1mm, colbacktitle=quarto-callout-note-color!10!white, leftrule=.75mm]

\begin{itemize}
\tightlist
\item
  \textbf{Sister chromatids} = identical copies (same alleles)
\item
  \textbf{Homologous chromosomes} = same genes, alleles may differ
\end{itemize}

\end{tcolorbox}

\begin{center}\rule{0.5\linewidth}{0.5pt}\end{center}

\subsection{Triploid, tetraploid,
polyploid}\label{triploid-tetraploid-polyploid}

\textbf{Diploid} = 2 complete sets\\
\textbf{Triploid} = 3 complete sets\\
\textbf{Tetraploid} = 4 complete sets\\
\textbf{Polyploid} = more than 2 sets :contentReference{oaicite:8}

\includegraphics[width=1\linewidth,height=\textheight,keepaspectratio]{../images/ploidy/polyploid_series.png}
\emph{Figure 9. One chromosome type shown across haploid → tetraploid.}

\begin{tcolorbox}[enhanced jigsaw, titlerule=0mm, coltitle=black, bottomrule=.15mm, colframe=quarto-callout-note-color-frame, title=\textcolor{quarto-callout-note-color}{\faInfo}\hspace{0.5em}{Why you'll see this later}, arc=.35mm, opacityback=0, opacitybacktitle=0.6, colback=white, breakable, toprule=.15mm, rightrule=.15mm, bottomtitle=1mm, left=2mm, toptitle=1mm, colbacktitle=quarto-callout-note-color!10!white, leftrule=.75mm]

Many plants are polyploid, and polyploidy is connected to traits like
larger fruit size. :contentReference{oaicite:9}

\end{tcolorbox}

\begin{center}\rule{0.5\linewidth}{0.5pt}\end{center}

\subsection{Practice: determine ploidy from a
picture}\label{practice-determine-ploidy-from-a-picture}

For each panel, answer: 1) What is \textbf{n}? (how many chromosome
types?)\\
2) How many complete sets are present? (ploidy)\\
3) What is the total chromosome number?

\includegraphics[width=1\linewidth,height=\textheight,keepaspectratio]{../images/ploidy/practice_panels.png}
\emph{Figure 10. Practice panels A--D.}

\begin{center}\rule{0.5\linewidth}{0.5pt}\end{center}

\subsection{Summary (memorize this structure, not
numbers)}\label{summary-memorize-this-structure-not-numbers}

\begin{itemize}
\tightlist
\item
  \textbf{n} = number of chromosome types in one complete set\\
\item
  \textbf{Ploidy} = number of complete sets of homologs\\
\item
  \textbf{2n} is diploid (two sets), \textbf{n} is haploid (one set)
  :contentReference{oaicite:10}
\item
  After S-phase, chromosomes are duplicated, but \textbf{ploidy does not
  change} (you still have the same number of sets)
  :contentReference{oaicite:11}
\end{itemize}

\begin{tcolorbox}[enhanced jigsaw, titlerule=0mm, coltitle=black, bottomrule=.15mm, colframe=quarto-callout-important-color-frame, title=\textcolor{quarto-callout-important-color}{\faExclamation}\hspace{0.5em}{One sentence to keep you honest}, arc=.35mm, opacityback=0, opacitybacktitle=0.6, colback=white, breakable, toprule=.15mm, rightrule=.15mm, bottomtitle=1mm, left=2mm, toptitle=1mm, colbacktitle=quarto-callout-important-color!10!white, leftrule=.75mm]

\textbf{Replication makes sister chromatids; it does not change ploidy.}

\end{tcolorbox}




\end{document}
